\section{Short products and the return to the accessible scalar}
\label{sec:short-products}

The scalar vacuum wave $l$ has not yet been realized as a physical field.
In this Section, $L$ is an insertion symbol whose one-vacuum wave is $l$;
it is not an operator on the polynomial field domain.  We construct only
products of total length at most four, including derivatives, and prove
exactly the factorization statements needed below.  The symbols
\begin{equation}
 J=V-\partial L,\qquad
 t=T-\frac13(\eta\Box-\partial\partial)L,\qquad
 \ell=L-S
 \label{P:formal-fields}
\end{equation}
initially have this same meaning.  At the end we prove $l=S$ as vacuum
waves.  Construction of mixed products at every finite length and of the physical
operator $L$ is deferred to Section~\ref{sec:reconstruction}.

We use the ordered-vector-tube convention of
Section~\ref{sec:extraction}: the vacuum is on the right, and
\begin{equation}
 \Im z_1\in V_+,\qquad \Im(z_{j+1}-z_j)\in V_+.
 \label{P:vector-tube}
\end{equation}
The strict first inequality is necessary for vector products.  Scalar
vacuum functions depend only on differences and require only the remaining
inequalities.  A derivative always acts on its own insertion.  All entries used in this Section are even.  Tensor
indices and source smearings in the following constructions are retained,
even when suppressed in the notation.

The length bound is chosen to support a specific Hilbert-space
argument.  Actual bounded products and finite interpolation will
first construct the short vectors and their split identities.  Their
four-point scalar functions then determine inner products of pairs,
which give a conformal representation on the pair space.  Its radial
positivity and the normalized current fluxes will force the residual
wave $l-S$ to vanish.

\subsection{Finite interpolation on spacelike patches}

\Needspace{6\baselineskip}
\begin{lemma}[Common wedge for small null flows]
\label{P:common-wedge}
Fix any finite family of sufficiently small, mutually spacelike real
source patches.  There are a point $o$ strictly to their future, an open set of
future null directions $n$, and, for each finite set of sufficiently small
null special-conformal parameters in that set, an opposite open cone
$O_-$ with the following property.  Bounded operations on the source
patches, or on their inverse coordinate images, after conjugation by the
actual centered null-charge unitaries lie in wedges spacelike to $O_-$.
The vectors $C\Omega$, with $C\in\mathcal M(O_-)$, form a total set.
These choices persist under sufficiently small changes of the source
patches and Lorentz transformations.
\end{lemma}

\begin{proof}
Choose a future null direction $m$ separated from the directions $n$ and
use $m$ as the conjugate wedge axis.  Work strictly below the null cut
plane centered at $o$.  In the resulting coordinates put $p(y)=|y|^2$.
Each source patch fits, with a strict margin, inside a cut
$W_{\alpha p+c}$ with $\alpha>0$ sufficiently small and $c$ sufficiently
negative.  The same containment holds for the finitely many inverse
coordinate images when the flow parameters are small enough.

Here is a two-sign bound uniform in quadratic truncation.  Choose smooth
bounded profiles $0\leq q_R\leq p$, increasing to $p$, and write
$U_R(t)=e^{itH_{q_R}}$.  For $0<t<\alpha$, cut inclusion and the
constant-shift identity give
\[
 \cM(W_{\alpha p+c})\subset\cM(W_{tq_R+c})
       =\Ad(U_R(t))\cM(W_c).
\]
Thus conjugation by $U_R(-t)$ sends the source algebra into the fixed
wedge algebra $\cM(W_c)$.  For positive $t$, the source algebra is
already contained in $\cM(W_c)$, and
\[
 \Ad(U_R(t))\cM(W_c)=\cM(W_{tq_R+c})\subset\cM(W_c).
\]
Both signs therefore have the same containing wedge for
$|t|\leq t_0<\alpha$, independently of $R$.  This argument uses
$q_R\leq p$, not a bound on the diverging norms $\|q_R\|_\infty$.
The positive moment construction of Section~\ref{sec:null-current}
gives $U_R(t)\to e^{itK_n}$ strongly.  Weak closedness of the fixed
wedge algebra now preserves the containment for each bounded conjugated
operation in this limit.  Translations give the centered version used
here.

For finitely many choices, take $m\cdot x$ sufficiently negative, then
add a sufficiently large positive multiple of $m$.  This leaves
$m\cdot x$ unchanged and makes all the finitely many $n\cdot x$
sufficiently positive.  A small open double cone about such an $x$ lies
in the required opposite region.  All inequalities are strict, hence
stable under the small changes asserted in the statement.  Local
vacuum cyclicity in this double cone proves totality.  No cyclicity of a
relative commutant is used.
\end{proof}

Let $g$ be the standard null special-conformal coordinate map centered at
$o$, with small parameter vector $b_g=sn$, and let $U_g$ denote the actual
unitary generated by the corresponding centered $K_n$.  On the
one-field space these actions agree canonically by
Theorem~\ref{C:canonical-Pi}.  This does not assert covariance on the
full physical Hilbert space.  Reparametrize an input at $g^{-1}x$ and
multiply its test and indices by the inverse primary factors at $x$.
The transformed one-vacuum data from accessible inputs $S,V,T$ are then
\begin{equation}
 S(x),\qquad
 V(e_j;x)+\frac{4b_g\cdot e_j}{1+2b_g\cdot(x-o)}L(x),\qquad
 t(x)+\mathcal P_g L(x).
 \label{P:transformed-data}
\end{equation}
Here $e_j$ is a fixed target polarization, and $\mathcal P_g$ is a
tensor-valued differential operator of order at most two with analytic
coefficients on the patch.  It is obtained by transforming
$(\eta\Box-\partial\partial)L/3$ as derivatives of a scalar primary of
weight two.  In particular
$\mathcal P_0=(\eta\Box-\partial\partial)/3$.
The coefficient in the vector formula is the derivative of the scalar
primary factor $(1+2b_g\cdot(x-o))^2$; this fixes its normalization.
All these identities also hold after differentiation and source
smearing.

\Needspace{6\baselineskip}
\begin{lemma}[Finite null interpolation]
\label{P:interpolation}
For $q\leq4$ and generic independent polarizations $e_1,\ldots,e_q$,
finite evaluations at arbitrarily small future null $b$ distinguish the
$2^q$ functions
\begin{equation}
 b\longmapsto\prod_{j\in I}
       \frac{4b\cdot e_j}{1+2b\cdot(x_j-o)},
 \qquad I\subset\{1,\ldots,q\},
 \label{P:interpolation-functions}
\end{equation}
uniformly for $x_j$ in sufficiently small compact source patches.
The inverse matrix is analytic in the positions, also in a complex
neighborhood where its determinant stays nonzero.
\end{lemma}

\begin{proof}
Complete the polarizations to a generic basis and put $z_j=b\cdot e_j$.
A polynomial in the square-free monomials $z_I$ that vanishes on an open
part of the null cone vanishes on the full complex quadric and is
divisible by its nondegenerate quadratic equation $Q(z)$.  This last
assertion follows, for example, by solving the quadric for one variable
on a nonsingular chart and applying polynomial division; the remainder
vanishes on an open set and is zero.  In a generic basis $Q$ has a nonzero
$z_j^2$ coefficient for every $j$.  If $Q R$ were nonzero and
multi-affine, then
$\deg_{z_j}(QR)=2+\deg_{z_j}R\geq2$, a contradiction.  Thus the
square-free monomials are independent on every such cone patch.
There is consequently a finite nonsingular evaluation matrix.  Varying
ray lengths separates its homogeneous degrees.

Scale all its sample vectors by $\varepsilon>0$ and divide column $I$
by $(4\varepsilon)^{|I|}$.  As $\varepsilon\downarrow0$, the matrix in
\eqref{P:interpolation-functions} converges uniformly on compact source
patches to the nonsingular monomial matrix.  Its determinant remains
nonzero for sufficiently small $\varepsilon$.  Matrix inversion gives
the last assertion.  Thus the interpolation involves multiplication of
distributions by smooth functions, never pointwise inversion of a
distribution.
\end{proof}

The limit $\varepsilon\downarrow0$ in the preceding proof is used
only to find a nonsingular evaluation matrix.  In the construction
below we fix one sufficiently small positive scale for the compact
placement or continuation geometry under consideration.  We then
shrink the position neighborhoods so that its determinant stays
bounded away from zero, and only afterward take the bounded-
approximation and boundary limits.  The inverse constants may depend
on this fixed scale; a uniform inverse as the scale tends to zero is
not needed.

Generic polarizations also suffice for exact tensor recovery.  The
decomposable tensors coming from admissible generic polarization
tuples span the full tensor product: a linear functional vanishing on
that open set is a multilinear polynomial vanishing identically.
Choose a finite spanning subset once.  Every requested tensor
component is then a fixed linear combination of those evaluations,
without inverting an evaluation matrix at a nongeneric tuple.

\Needspace{6\baselineskip}
\begin{proposition}[Hilbert products through length four]
\label{P:short-products}
For entries from $t,S,V,L$, and hence $J,\ell$ and their derivatives,
there are Hilbert-valued product distributions on mutually spacelike
real patches through total length four.  They are multilinear,
permutation-symmetric, and covariant under translations, Lorentz
transformations and dilations.  They satisfy every leg-local linear
identity of their vacuum waves.  With $C$ in a common opposite cone,
they obey
\begin{equation}
 \ip{C\Omega}{\Psi[O_1\cdots O_p]}
 =\ip{O_1^\dagger\Omega}{C^*\Psi[O_2\cdots O_p]}.
 \label{P:recursive-product}
\end{equation}
Smearings on the left side of the inner product are conjugated.
The scalar functions factor at every split of total length at most
four.  Explicitly, for $p+q\leq4$ and both words in their vector tubes,
\begin{equation}\label{P:short-splitting}
 \begin{split}
 &\big\langle\Psi[A_1(z_1)\cdots A_p(z_p)],
                  \Psi[B_1(w_1)\cdots B_q(w_q)]\big\rangle\\
 &\quad=W\bigl(A_p^\dagger(\bar z_p),\ldots,
               A_1^\dagger(\bar z_1),B_1(w_1),\ldots,B_q(w_q)\bigr).
 \end{split}
\end{equation}
Here an empty word is $\Omega$, and bra source tests are conjugated
and reversed.  The reversed bra gaps, the ket gaps, and the central
gap $\operatorname{Im}w_1+\operatorname{Im}z_1$ are all future
directed.  Thus the scalar word is in its gap-only tube.  In particular
the case $p=q=2$ gives the pair--pair Gram identities.  All these
products extend holomorphically to
\eqref{P:vector-tube} and the spacelike tie neighborhoods of
Proposition~\ref{E:accessible-products}; the split identities persist
wherever both vector factors are defined.
\end{proposition}

\begin{proof}
Proceed by induction on length, and at a fixed length by increasing
number of $t$ slots.  Length one is the given one-vacuum construction.
For accessible inputs use the localized bounded approximants of
Proposition~\ref{E:bounded-approximation}.  Conjugate every factor of an
approximating real product by the same $U_g$.  Lemma~\ref{P:common-wedge}
allows the first factor to commute past $C$.  Vacuum-vector and
adjoint-vacuum convergence for this factor, and convergence of the
shorter tail, therefore yield the right side of
\eqref{P:recursive-product}.  By the shorter induction assertions this
is the multilinear expansion of \eqref{P:transformed-data}, initially
as a functional of $C\Omega$.

Use $T$ in positions intended eventually to contain $t$.  Terms with
fewer $t$ slots have representing vectors already and can be subtracted:
replacing a transformed $T$ by $\mathcal P_gL$ lowers this induction
index.  Interpolation is then needed only for the remaining $V/L$
choices.
At the remaining $V$ positions the unknown functionals are indexed by
the subsets replaced by $L$, with the coefficients
\eqref{P:interpolation-functions}.  Lemma~\ref{P:interpolation} expresses
each unknown as a finite smooth linear combination of already existing
Hilbert-valued distributions.  They therefore have representing
vectors.  Generic polarization choices suffice: multilinearity,
followed by a finite choice of generic bases, recovers all components.
This proves \eqref{P:recursive-product} at the new length.

The same argument with any other sufficiently small $g$ proves the
expansion formula, because it proves equality against the total family
$C\Omega$ in a common opposite cone.  Subdivision of patches or a
change in all auxiliary choices gives the same vectors recursively by
\eqref{P:recursive-product}.  Comparing a construction with a small
translated, Lorentz-transformed or rescaled construction proves the
similarity laws.  Interpolating the symmetry of transformed accessible
products proves spacelike permutation symmetry, again first subtracting
terms with fewer $t$ slots.  Inserting a zero vacuum wave into
\eqref{P:recursive-product}, after permuting it to the first position,
gives zero.  Hence, in particular,
\begin{equation}
 \Box L=\partial\cdot V=\Theta,\qquad
 \partial\cdot J=0,\qquad
 \partial^\mu t_{\mu\nu}=0,\qquad t^\mu{}_{\mu}=0
 \label{P:leg-identities}
\end{equation}
hold in these products.  This argument concerns vector distributions,
not operator-domain equalities.

For bounded approximants the vacuum expectation of a product splits at
any position by Hilbert-space adjunction.  The same is true for its
common unitary conjugate.  Expand the full scalar product and both
factors by \eqref{P:transformed-data}; real transformations respect
adjunction.  Subtract the terms already split by induction in the
number of $t$ slots.  The identical finite interpolation now proves
\eqref{P:short-splitting} for every remaining entry list, with the
same adjoint and coordinate convention as
\eqref{E:accessible-splitting}.  Thus short-length
associativity and positivity come from actual bounded products.

It remains to continue the vectors, rather than merely their matrix
coefficients.  Fix a compact continuation path from a spacelike
boundary patch to an interior tube point or one of the allowed ties.
Choose the parameter scale in Lemma~\ref{P:interpolation} small enough
that every inverse image $g^{-1}z_j$ stays in the accessible product
domain along that path.  Near a tie use its open holomorphic
neighborhood.  The analytic inverse matrices remain nonsingular on a
small neighborhood of the compact path.  Applying $U_g$ to the
accessible analytic vectors and performing the same successive
subtractions and inversions constructs the desired holomorphic
vectors.  The boundary calculus of
Lemma~\ref{E:boundary-calculus} identifies their distributional
boundary with the real construction: the coordinate maps take a
smaller approach wedge strictly into the original wedge.

There is no choice of branch.  On overlapping connected compact
continuation neighborhoods choose parameters sufficiently small for
both constructions; their boundaries agree on a smaller original
patch, and distributional edge-of-the-wedge uniqueness makes the
holomorphic vectors agree.  Exhaustion gives the full ordered tube.
At a tie the accessible germs and previously constructed terms agree
across the relevant permutations, so the interpolated germs agree as
well.  All leg identities and similarity laws continue by the same
uniqueness.  Finally, the two sides of each split identity have equal
boundary distributions when bra and ket are separately smeared on the
original spacelike patches and approached in strict subcones.
Uniqueness in their common split tube proves the asserted continuation
of the identity.  Only a local distributional boundary and an open
convex approach subcone are needed here; no global boundary estimate
for arbitrary $L$ products has been assumed.
\end{proof}

\Needspace{6\baselineskip}
\begin{lemma}[Comparison with the single-insertion forms]
\label{P:single-comparison}
The analytic pairs $LV,VL,JV,VJ,tV,Vt$ agree, when paired with an
ordinary physical polynomial bra, with the corresponding insertions of the
tempered single-$L$ forms of Proposition~\ref{C:single-insertion}.
\end{lemma}

\begin{proof}
First take a finite sum of factorized compact physical bra words and
translate it into a common opposite cone.  Its affiliated bounded
truncations can be used as $C$ in
\eqref{P:recursive-product}.  With $L$ first, the adjoints of these
truncations acting on an ordinary smeared $V\Omega$ converge on the
polynomial domain.  Locality of the single-insertion forms therefore
gives equality on an open real spacelike patch.  For a compact joint
bra kernel, approximate by factor tests in the permitted smooth
topology, preserving the opposite-cone margin.  Continuity of the
vacuum and adjoint-vacuum graph vectors and of the single-insertion
distributions extends this equality to that kernel.  Give bra and pair
strictly separated ordered imaginary translations.  The known
single-insertion spectral cones give analyticity in the two pair
arguments and the common translation of the bra, with its internal
spectator tests fixed.  Analytic uniqueness extends the equality, and
one may remove the separate bra damping while keeping the pair damped.
The other pairs follow from their linear differential definitions.
\end{proof}

\subsection{The Ward rows and the alternating exception}

The vector construction supplies finite transformation identities
for products formed from the accessible entries.  To obtain the
canonical Ward identity for each interpolated entry separately, we
must recover every row of the triangular system mixing $V$ and $L$.
The possible four-slot alternating term is the one case not removed
by the square-free interpolation argument; its elimination uses the
purely timelike spectrum of the last scalar wave.

\Needspace{6\baselineskip}
\begin{proposition}[Conformal covariance of the short scalar functions]
\label{P:short-covariance}
All scalar vacuum functions $W$ of total length at most four obey the
infinitesimal conformal Ward identities, with $L,S,J,t$ primary of
weights $2,2,3,4$ and spins $0,0,1,2$, respectively.  The laws for $V,T$
are their derivative laws from \eqref{P:formal-fields}.
\end{proposition}

\begin{proof}
Fix the $t$ and $S$ positions.  At the remaining $m_0\leq4$ positions
start with polarized vectors $V(e_j)$, and let $W_I$ be the function
obtained by replacing the positions in $I$ by $L$.  The canonical
special-conformal Ward operator $\widehat K_a$ acts diagonally by the
primary differential actions and has off-diagonal part
\begin{equation}
 (\widehat K_a W)_I\big|_{\mathrm{off}}
       =4\sum_{j\notin I}(a\cdot e_j)W_{I\cup\{j\}}.
 \label{P:ward-off-diagonal}
\end{equation}
These operators commute: they are the canonical commuting special
transformations of primaries and their derivatives.  They may be
centered at $o$, since the difference from another center is an affine
Ward operator, already zero on every row.

Unitarity and the expansion in Proposition~\ref{P:short-products} give
\begin{equation}
 (e^{s\widehat K_n}W)_\varnothing=W_\varnothing
 \label{P:null-row}
\end{equation}
for small $s$ and the open family of null $n$.  Initially the tensor
positions contain $T$.  Induct in their number, substitute
$T=t+\mathcal P_0L$, and subtract terms with fewer $t$ positions, whose
covariance is already established.  This gives
\eqref{P:null-row} with the indicated $t$ positions.

Write $X_I(a)=(\widehat K_aW)_I$.  Differentiating
\eqref{P:null-row} once and using the linear span of the open null family
gives $X_\varnothing(a)=0$ for every $a$.  Suppose that $X_I=0$ when
$|I|<k$.  Commuting two Ward operators on a row $H$ of size $k-1$ and
differentiating \eqref{P:null-row} $k+1$ times give, respectively,
\begin{align}
 \sum_{j\notin H}z_j(a)X_{H\cup\{j\}}(b)
   &=\sum_{j\notin H}z_j(b)X_{H\cup\{j\}}(a),
       \label{P:row-symmetry}\\
 \sum_{|I|=k}z_I(n)X_I(n)&=0,
       \label{P:row-null}
\end{align}
where $z_j(a)=a\cdot e_j$ and $z_I=\prod_{j\in I}z_j$.
Indeed all diagonal differentiations of already zero rows vanish; the
first surviving paths through the triangular off-diagonal array are
exactly the $k!$ orders of adding the same $k$ indices.  The common
nonzero factors $4$ and $k!$ have been canceled.

Complete a generic set of polarizations to a basis and write
$X_I(a)=\sum_{h=1}^4z_h(a)x_{Ih}$.  Equation~\eqref{P:row-symmetry}
says that, for each $H$, the matrix with rows outside $H$ equal to
$x_{H\cup\{j\},h}$ and its other rows zero is symmetric.  If $k\geq2$,
choose $H=I\setminus\{j\}$.  For $h\in I$ choose $j\ne h$; its zero
row kills $x_{Ih}$.  A completed-basis row $h>m_0$ likewise kills that
entry.  For $h\notin I$, symmetry identifies all the surviving entries
with a coefficient $c_{I\cup\{h\}}$, depending only on the union.
Thus \eqref{P:row-null} becomes
\begin{equation}
 (k+1)\sum_{|B|=k+1}c_Bz_B(n)=0.
 \label{P:square-free-ward}
\end{equation}
The square-free independence proved in
Lemma~\ref{P:interpolation} kills every coefficient, including the case
where there is no such union.

For $k=1$, \eqref{P:row-symmetry} instead gives a symmetric matrix.
Its quadratic form vanishes on the null cone, so it equals a scalar
multiple of the metric quadratic.  If $m_0<4$, its zero completed-basis
rows force that multiple to vanish.  If $m_0=4$, let
$G_{ij}=e_i\cdot e_j$.  The coefficient matrix is $\beta G^{-1}$,
which says
\begin{equation}
 X_{\{i\}}(e_j)=\beta\delta_{ij}.
 \label{P:exception-evaluation}
\end{equation}
The function $\beta$ is multilinear in all four polarizations: row $i$
is independent of $e_i$, and evaluation on $e_i$ restores its linear
factor.  It vanishes when two polarizations coincide, by
\eqref{P:exception-evaluation} and continuity from generic choices.
It is therefore alternating.  The sole possible exception is
\begin{equation}
 X_{\{i\}}(a)=
 \lambda\,\epsilon(e_1,\ldots,e_{i-1},a,e_{i+1},\ldots,e_4).
 \label{P:alternating-exception}
\end{equation}
Here there are exactly four $V/L$ positions and no other entries.
Divergence in position $i$ of $X_\varnothing=0$ also gives
\begin{equation}
 \Box_iX_{\{i\}}=0.
 \label{P:exception-harmonic}
\end{equation}
For clarity, the canonical array preserves the relation
$\partial\cdot V=\Box L$: its difference is the divergence of the
primary dimension-three vector $J$, whose divergence transforms
homogeneously.  This is the intertwining identity used in deriving
\eqref{P:exception-harmonic}.

We eliminate \eqref{P:alternating-exception} using its last position,
$i=4$.  Smear the first three positions on a real mutually spacelike
patch, keeping the last argument in the future tube.  Split the scalar
function into the triple bra and the one-wave ket.  Its boundary in the
first three positions is obtained from negative ordered bra dampings.
Move the center of the generator sufficiently far to the future of
these compact supports; affine covariance leaves $X$ unchanged.
The last-position canonical differential action on $L$ is exactly the
canonical generator applied to $l(z_4)$.  The ordinary physical triple
$VVV$ bra belongs to the domains of these centered actual $K_n$ by
Lemma~\ref{C:K-domains}.  Moreover $\Pi$ reduces these
generators and their actions on $\Pi$ are canonical by
Theorem~\ref{C:canonical-Pi}.  Projecting the bra to $\Pi$ thus
allows the last-position generator to be moved to the fixed bra.
Linear combinations of the available null parameters span arbitrary
$a$.

Every other term of the smeared $X_{\{4\}}(a)$ pairs the unchanged
last wave with a fixed constructed triple bra, possibly differentiated
or with replacements.  Consequently the whole expression is
$\ip{\eta}{l(z_4)}$ for a fixed Hilbert vector $\eta$, with no
coordinate multiplier left on the last wave.  Its translation
spectrum is purely timelike, since that of $l$ is.  Equation
\eqref{P:exception-harmonic} says that this spectrum is supported on
$p^2=0$; these disjoint supports force it to vanish.  More explicitly,
after any fixed future damping it is a finite spectral pairing, and
multiplication by $p^2$ kills it only if that pairing is zero off the
null shell.  There is no shell or zero-momentum component of $l$.
Boundary uniqueness in the first three arguments and generic
polarization then give $\lambda=0$.

The induction on $k$ is now complete, followed by the induction in the
number of $t$ positions.  This proves every row of the canonical Ward
system.
\end{proof}

\subsection{Euclidean continuation and the residual at infinity}

Wick coordinates are $(\tau,\boldsymbol r)\mapsto i\tau e+\boldsymbol r$.
Sort points by increasing $\tau$ and use the tie neighborhoods when
heights coincide.  A distinct Euclidean configuration has only
spacelike equal-height ties, so the resulting scalar functions are
real analytic on the full distinct-point configuration space.  Use
complex-linear tensor pullback with basis $C_0=ie$ and real spatial
bases $C_j$.  The identities become
\begin{equation}
 V=J+\partial L,\qquad
 T=t+\frac13(\delta\partial^2-\partial\partial)L.
 \label{P:euclidean-identities}
\end{equation}
Here $\partial^2$ is the Euclidean Laplacian.  Integrating the Ward
identities along nonsingular local conformal flows gives their finite
covariance, with reordering through tie neighborhoods as required.

For a scalar primary $A$ of dimension two define
\begin{equation}
 W(A_\infty\,\cdots)=\lim_{R\to\infty}R^4W(A(Ru)\,\cdots),
 \qquad |u|=1.
 \label{P:infinity-definition}
\end{equation}
The limit exists and has the primary transport law.  To verify this,
choose a Euclidean special-conformal transformation carrying infinity
to a finite point disjoint from the spectators.  Join it to the
identity by a flow whose pole ray avoids the finite set of points,
including the moving large point.  At the final transformation every
argument is distinct and finite, so ordinary local analyticity and the
weight-two factor give the limit.  This argument is uniform on small
complex neighborhoods of strictly time-ordered spectator
configurations and proves locally holomorphic convergence there.  It
also proves independence of the direction in
\eqref{P:infinity-definition}.

\Needspace{6\baselineskip}
\begin{proposition}[Residual infinity decoupling]
\label{P:residual-infinity}
For the residual scalar $\ell=L-S$,
\begin{equation}
 W(\ell_\infty O_1\cdots O_j)=0,
 \qquad O_i\in\{S,V,T\},\qquad j\leq3,
 \label{P:accessible-decoupling}
\end{equation}
at all distinct Euclidean points, also after differentiating any
spectator insertion.
\end{proposition}

\begin{proof}
We justify passage from the localized-vector result of
Corollary~\ref{E:infinity-vanishing} to these products; unrestricted
timelike real boundaries for accessible products are unnecessary.
For a real vector localized in a fixed line diamond, the projected
Paley multiplier in that Proposition has zero value at energy zero.
The same local Paley bound controls its first derivative there.  Thus
it is $O(E)\norm{\xi}$ at low energy.  Integrating with the
weight-two energy powers against $e^{-RE}$ in the measure $\dd E/E$,
and using exponential damping beyond a fixed energy, gives
\begin{equation}
 R^4\abs{\ip{\ell(iRe)}{\xi}}
       \leq C R^{-1}\norm{\xi},\qquad R\geq1.
 \label{P:infinity-rate}
\end{equation}
For complex localized operator data the right side is bounded by the
sum of the vacuum and adjoint-vacuum norms.  Consequently the estimate
passes to mutually spacelike accessible products in that diamond by
the two bounded-approximation limits of
Proposition~\ref{E:accessible-products}.

Choose tests on small patches about equal-time mutually spacelike
points and continue their translations a short distance into strictly
increasing imaginary time.  Before using the improved real-edge
estimate, the corresponding smeared pairing multiplied by $R^4$ is
$O(R^2)$ uniformly up to a smaller good real patch: the one-wave norm
is $O(R^{-2})$ and the smeared product bounds are uniform there.  Apply
the subharmonic logarithmic estimate on a half-disk with this real
edge.  If its harmonic measure at an interior point is $\omega$, the
bound becomes
$O(R^{-\omega+2(1-\omega)})=O(R^{2-3\omega})$.
There is an interior open set with $\omega>2/3$, and hence convergence
to zero on that set.  The positive imaginary times can be fixed before
further shrinking the source tests.  On its strict Euclidean
configurations \eqref{P:infinity-definition} already converges locally
holomorphically.  Small real source offsets therefore commute with
that limit.  Unsmeared uniqueness gives
\eqref{P:accessible-decoupling} on an open set, and real analyticity
continues it to every distinct configuration.  Derivatives follow
from the same locally holomorphic identity.
\end{proof}

\subsection{The pair representation and normalized fluxes}

The four-point split identities now supply a Hilbert-space conformal
action on pairs. Its normalized current fluxes will control the possible
singular terms in the radial argument below.

\Needspace{6\baselineskip}
\begin{proposition}[Conformal representation on the pair space]
\label{P:pair-representation}
Let $\mathcal H_{\leq2}$ be the closed span of $\Omega$ and the one- and
two-insertion tube vectors for $L,S,J,t$.  There is a strongly
continuous unitary representation of the universal cover of the
connected conformal group on this space, with the canonical local
primary transformations.  Translations coincide with the actual
positive-energy translations.  Its timelike special-conformal
generator $K_e^{\mathrm{rep}}$ is nonnegative.
\end{proposition}

\begin{proof}
For each length restrict its arguments to a relatively compact open
box in its tube.  These vectors span a dense subspace: a vector
orthogonal to them gives a holomorphic scalar function zero on an
open set, and hence zero throughout the tube.  For all sufficiently
small real conformal transformations, map these vectors by their
primary transformation laws.  The resulting inner products are
unchanged by the split identities of
Proposition~\ref{P:short-products} and the integrated scalar Ward
identities of Proposition~\ref{P:short-covariance}.  Total length four
is exactly what is needed for pair--pair inner products.  The image
is dense by the same open-box argument; thus the map extends to a
unitary.  On smaller boxes the local composition identity is exact.
Holomorphic point dependence gives strong continuity.  Continuing
these local unitaries along paths integrates the action to the
universal cover.

Uniqueness extends the local laws to all tube points and, for
sufficiently small transformations, to Euclidean configurations in
the positive-time half-space.  The point vectors there are smooth
group vectors.  On translations the construction agrees with the
actual translation action, whose spectrum is future directed.  The
time-M\"obius subgroup conjugates its positive time-translation
generator to $K_e^{\mathrm{rep}}$, proving the final assertion.  No
conformal action on the full physical Hilbert space is asserted here.
\end{proof}

For a target $A(y)$ in the positive Euclidean time half-space let
$\Sigma$ be a small outward-oriented sphere about $y$, wholly within
that half-space.  Define Hilbert-valued fluxes from the analytic pair
vectors by
\begin{equation}
 F_J(A;y)=\int_\Sigma\dd\Sigma^a\,\Psi[J_a(x)A(y)],\qquad
 F_X^t(A;y)=\int_\Sigma\dd\Sigma^a\,X^b(x)\Psi[t_{ab}(x)A(y)].
 \label{P:flux-definition}
\end{equation}
At each point of the sphere use the Euclidean order, with the tie
germs providing the transitions.  There is no factor of the sphere
area in \eqref{P:flux-definition}.  Conservation and tracelessness
make the corresponding holomorphic flux forms closed off the target,
so the contour may be deformed, or held fixed when differentiating
the target nearby.  If one instead rescales both currents' fluxes by
a common conventional constant, the constant $c$ below rescales with
them.

\Needspace{6\baselineskip}
\begin{proposition}[Short-length Ward fluxes]
\label{P:flux-wards}
With one consistent nonzero Ward normalization $c$,
\begin{align}
 F_J(A)&=0,& A&=V,L,J,\label{P:J-flux}\\
 F^t_{\partial_b}(A)&=c\,\partial_b A\Omega,& A&=V,L,J.
 \label{P:translation-flux}
\end{align}
For primary targets $L,J$ and every conformal Killing field $X$,
\begin{equation}
 F_X^t(A)=c\,\delta_X A\Omega.
 \label{P:conformal-flux}
\end{equation}
Here $\delta_X$ is the canonical differential action, including spin
and weight.  In particular, a special transformation centered at the
target has zero flux, while centered dilation about $L$ gives
$2cL\Omega$.  With outward orientation in
\eqref{P:flux-definition}, Wick time basis $C_0=ie$, the convention
$i[P_\nu,A]=\partial_\nu A$, and upper-minus-lower order $At-tA$ in
the cap difference, one has $c=-1$.
\end{proposition}

\begin{proof}
First take $A=V$ and pair with an arbitrary ordinary physical
polynomial bra.  Lemma~\ref{P:single-comparison} identifies the
analytic pairs with tempered single-insertion forms.  Deform
$\Sigma$ to a cylinder whose relative Euclidean time caps are
$\pm\delta$ and whose spatial disk has radius $a$.  Give the whole
relative time a small real Minkowski shift $s$, with $|s|<a$.
The shifted contour remains in the holomorphic domains; where the
imaginary relative time crosses zero on the side, strict spatial
separation and tie gluing provide a common neighborhood.

Average smoothly and with integral one over $a$ and $s$, keeping the
shift support strictly inside $|s|<a$.  The side contribution tends to
zero as $\delta\downarrow0$ because it lies on a compact nonsingular
spacelike set and its time thickness tends to zero.  The two caps now
have smooth Cauchy-slab tests and converge to the opposite Lorentzian
order boundaries at a common target damping.  This use of the
single-insertion boundary is legitimate even at a point in the common
complex target translation: its total pair momentum is exponentially
damped, while the smooth relative test decays rapidly in the relative
momentum.  The difference is the commutator charge with the time
average and a spatial cutoff equal to one on its causal support.
The $J$ charge is zero by Proposition~\ref{C:Q-zero}; the translation
charge of $t$ is the original translation charge by
Section~\ref{sec:charges-cones}.  They give
\eqref{P:J-flux}--\eqref{P:translation-flux} for $V$.

The Wick time component contributes $i$, and the tensor translation
index contributes the pullback $C_b^\nu$.  With the order and charge
convention in the statement,
$i[V,P_b]=-\partial_bV$, hence $c=-1$.  This calculation fixes one
common normalization for all components and all subsequently
transported Killing fields.  Pairing with the dense polynomial bras
proves equality of the Hilbert vectors.

A constant-current flux around a primary target extends to a
holomorphic single-wave function of its Minkowski argument throughout
the future tube: choose a relative sphere smaller than its common
damping.  Translation covariance makes it a translation wave;
dilations give its flux dimension.  Small Lorentz transformations act
on the pair and its contour; homotopy of the transformed contour back
to a small sphere gives the tensor transformation.  Differentiate the
identities for $V$ in its target index and use
$\partial\cdot V=\Box L$ while keeping the enclosing contour fixed.
One obtains the homogeneous wave equation for both
\begin{equation}
 F_J(L),\qquad
 F^t_{\partial_b}(L)-c\,\partial_bL\Omega.
 \label{P:harmonic-flux-defects}
\end{equation}
These are, respectively, a scalar wave of dimension two and a vector
wave of dimension three.  Translation covariance and positive-energy
holomorphy realize their values after any fixed future damping as
vacuum spectral amplitudes; further damping agrees by uniqueness.
Their wave equation confines their spectrum to $p^2=0$.
Scale covariance excludes a zero-momentum amplitude, and the
light-cone boost-weight test of Section~\ref{sec:foundations} excludes
the nonzero shell: in these cases $d-1>\max\omega$ for the finite
Lorentz multiplet.  Only local finiteness after exponential damping is
needed for that test.  Thus both defects vanish.  Subtraction of
$\partial L$ from $V$ proves the identities also for $J$.

For $A=L,J$, conjugate the translation identity by any sufficiently
small real conformal transformation in
Proposition~\ref{P:pair-representation}.  Since $t$ is a conserved
traceless primary of dimension four, its flux transforms exactly to
the flux for the pushed-forward vector field: three surface powers,
one vector-field power, and the two orthogonal index factors cancel
its weight.  Deform the transformed sphere back locally.  The target
primary factors are invertible, and the differential action on the
right transforms by the same Lie algebra adjoint action.  Translations
and all their such adjoints span the conformal Lie algebra (successive
conformal brackets give rotations, dilations and special
transformations).  Linearity and complexification give every
Euclidean conformal Killing field and prove
\eqref{P:conformal-flux}.
\end{proof}

The flux laws remain valid inside scalar Euclidean functions up to the
available split length, with exterior spectators, including
$\ell_\infty$.  First place those spectators earlier than the whole
pair by a Euclidean time plane and use the split identity, representing
infinity first by a large negative-time insertion.  Then continue in
the distinct spectator and target positions, choosing a smaller
contour where necessary.  One immediate consequence is
\begin{equation}
 W(\ell_\infty J_a(x)L(y))=0.
 \label{P:ell-J-L}
\end{equation}
Indeed rotations, translations and dilations give the only possible
form $d(x-y)_a/|x-y|^4$.  Its $J$ flux about $L$ is $d|S^3|$, and
\eqref{P:J-flux} makes $d=0$.  Primary covariance transports this
vanishing to any placement of the infinity insertion.

We now have the pair representation and current flux identities
with one fixed normalization.  The remaining argument turns radial
singularities of short functions into spectral statements on that
pair space.  Only the representation classifications at energies
zero, one and two will be needed for this step, so the analysis does
not require a discrete expansion of
all radial states.

\subsection{Radial Hilbert calculus and its low energies}

The function to be tested is
$W(\ell_\infty\ell(x)V_\mu(y)V_\nu(z))$.
Residual infinity decoupling will make it harmonic in $x$.  We will
identify each singular radial power with an actual low-energy spectral
projection on the pair space.  The low-energy bounds and a Maxwell
test leave only dipoles; a tangential reflection then removes those
dipoles from a positive pair norm.  The resulting norm bound eliminates
the remaining low-energy overlaps, and the normalized current fluxes
will finally test the residual scalar's own norm.  The following
calculus supplies these spectral and positivity statements without a
complete expansion into radial eigenstates.

Write a Euclidean point as $y=(y_0,\boldsymbol y)$.  The conformal map
\begin{equation}
 c_{\mathrm{ball}}(y)=
 \frac{(1-|y|^2,\,2\boldsymbol y)}{1-2y_0+|y|^2}
 \label{P:ball-map}
\end{equation}
takes the unit ball to the positive-time half-space.  Define radial
one- and two-insertion kets by primary pullback of the corresponding
half-space vectors; obtain descendants by differentiation.  Radial
adjunction reflects $y$ to $Iy=y/|y|^2$ with its full tensor and
weight factors.  Indeed direct substitution in
\eqref{P:ball-map} gives
$c_{\mathrm{ball}}(Iy)=(-c_{\mathrm{ball}}(y)_0,
 c_{\mathrm{ball}}(y)_{\mathrm{spatial}})$, exactly Wick reflection.
Scalar covariance therefore identifies radial splitting with Hilbert
adjunction.  To transport arbitrary separated interior and exterior
groups, each of size at most two, join the required proper Euclidean
conformal transformations by paths avoiding poles on the finite point
set.  A pole condition has codimension four, so such paths can be
chosen after a small perturbation.  Points at infinity are reached by
\eqref{P:infinity-definition}.  In particular the radial ket
$|\ell\rangle=|\ell(0)\rangle$ has bra $\ell_\infty$.

On the pair representation put
\begin{equation}
 H=\frac12(P_e+K_e^{\mathrm{rep}}).
 \label{P:radial-H}
\end{equation}
This is the self-adjoint compact-time generator on the covering group.
Its derived-generator identity holds on group-smooth vectors, a core
for this generator.  Both terms on the right are positive, so $H\geq0$.
The real Minkowski vector field of $H$, in Wick half-space coordinates,
is
$i((\tau^2-|\boldsymbol r|^2-1)/2,\tau\boldsymbol r)$.
The derivative of \eqref{P:ball-map} takes $iy$ to this field.  Thus
positivity and analytic continuation of the local group law give
\begin{equation}
 r^H|A_1(y_1)\cdots A_p(y_p)\rangle
   =r^{d_1+\cdots+d_p}
       |A_1(ry_1)\cdots A_p(ry_p)\rangle,
 \quad 0<r\leq1,\quad p\leq2.
 \label{P:radial-contraction}
\end{equation}
One first continues through a short contraction, where both analytic
laws apply, and then iterates through interior configurations.
Radial rotations act unitarily with the usual $SO(4)$ tensor law.

The real-form involution on complex radial transformations is complex
conjugation followed by conjugation with $I$.  Write $p_a^+,p_a^-$,
$H,M_{ab}$ for the derived actions corresponding, respectively, to
radial translation, $2y_a y-|y|^2e_a$, dilation, and
$y_a e_b-y_b e_a$.  On smooth vectors,
\begin{equation}
 (p_a^+)^\dagger=p_a^-,\qquad
 [H,p_a^\pm]=\pm p_a^\pm,\qquad
 [p_a^-,p_b^+]=2\delta_{ab}H+2M_{ab}.
 \label{P:radial-algebra}
\end{equation}
These are active-point conventions.  For example,
$M_{ab}|A_d\rangle=\delta_{ad}|A_b\rangle-\delta_{bd}|A_a\rangle$
at the origin.  Inversion takes translations to negative special
transformations, giving the displayed matrix adjunction.

The positivity calculation below is the low-energy part of the
four-dimensional conformal unitarity-bound method of
Mack~\cite{Mack1977}; see also the level-one Gram-matrix derivation
in \cite[Section~2.5]{Minwalla1998}.  We give the required finite-array
calculation on the present pair representation, including its
smooth-density justification.

\Needspace{6\baselineskip}
\begin{lemma}[The required low radial energies]
\label{P:radial-low-energy}
The $H=0$ eigenspace is a trivial representation.  The $H=1$ eigenspace
has only scalar rotation types.  At $H=2$, among the types in a tensor
product of two vectors, only scalars and antisymmetric two-forms can
occur.  Every smooth antisymmetric energy-two array obeys
\begin{equation}
 \sum_a p_a^+|F_{ab}\rangle=0,\qquad
 \sum_a p_a^+|(*F)_{ab}\rangle=0.
 \label{P:Maxwell-null-vectors}
\end{equation}
No discreteness of the full radial spectrum is required.
\end{lemma}

\begin{proof}
We specify the smooth-density issue first.  The element
$Z=e^{2\pi iH}$ is central, since its adjoint action on the conformal
Lie algebra is the identity by \eqref{P:radial-algebra}.  Its fixed
subrepresentation contains every integer $H$ eigenspace.  The
projection onto that subrepresentation commutes with the group,
preserves smooth vectors and has dense smooth image.  There the
spectrum of $H$ is supported on the integers, and
\begin{equation}
 E_H(\{h\})v=\frac1{2\pi}\int_0^{2\pi}e^{it(H-h)}v\,\dd t
 \label{P:integer-projection}
\end{equation}
preserves smoothness.  Averaging also over the compact rotations gives
dense smooth vectors in each rotation type and dense smooth
finite tensor-equivariant arrays.  Rotations descend to $SO(4)$,
since their action on the dense radial kets is the ordinary tensor
action.  The computations below therefore extend by density to the
stated eigenspaces.

At zero energy the expectations of both positive terms in
\eqref{P:radial-H} vanish.  They annihilate a smooth zero-energy vector.
The translation spectrum cone then makes every translation generator
annihilate it, and the conjugate cone makes every special partner
annihilate it.  Their brackets annihilate it as well, proving
triviality.  If $Hv=v$, lowering has energy zero and
$\ip{w}{p_a^-v}=\ip{p_a^+w}{v}=0$ for every zero-energy $w$;
hence $p_a^-v=0$.

For any smooth lowest array of energy $h$, raising norms give
\begin{equation}
 \left\|\sum_a p_a^+v_a\right\|^2
   =2\sum_{a,b}\ip{v_a}{(h\delta_{ab}+M_{ab})v_b}\geq0.
 \label{P:raising-Gram}
\end{equation}
A two-by-two plane-rotation block gives $|m|\leq h$ for each plane
weight.  An $SO(4)$ tensor type has two $SU(2)$ spins $j_+,j_-$, their
sum an integer, with maximal plane weight $j_++j_-$.  At $h=1$ only
scalars, vectors and the two chiral two-form types survive this test.
The divergence arrays exclude the last two.  In detail, the tensor
index rotation rule gives
\begin{align}
 \sum_bM_{ab}|A_b\rangle&=-3|A_a\rangle,\label{P:vector-Gram}\\
 \sum_cM_{ac}|F_{cb}\rangle
   &=-3|F_{ab}\rangle-|F_{ba}\rangle
                  +\delta_{ab}\sum_c|F_{cc}\rangle.
 \label{P:tensor-Gram}
\end{align}
The matrix eigenvalues in \eqref{P:raising-Gram} are therefore $-3$
for the vector divergence, $-4$ for a symmetric traceless rank-two
divergence, and $-2$ for an antisymmetric rank-two divergence.
Their positivity bounds are $h\geq3,4,2$, respectively.  At $h=1$
this leaves only scalars.

At $h=2$ the types in two vectors are even under the central radial
rotation $y\mapsto-y$.  Their lowered arrays are odd but would have to belong to
the scalar-only energy-one space, so they are lowest arrays.  The
$-4$ bound excludes the symmetric traceless type.  The antisymmetric
bound is saturated and makes the first norm in
\eqref{P:Maxwell-null-vectors} zero.  Apply the same argument to the
Hodge dual, which intertwines rotations, to obtain its second
identity.  This proves exactly the needed low-energy statements
without an expansion in a complete family of local operators.
\end{proof}

\Needspace{6\baselineskip}
\begin{lemma}[The exterior Maxwell test]
\label{P:Maxwell-exterior}
An antisymmetric energy-two radial array has zero exterior overlap
with the pair $\ell_\infty,V_\nu(z)$, with $z$ outside the radial ball.
\end{lemma}

\begin{proof}
Let $B_\nu(z)$ be the reflected ket of this exterior data; it is a
smooth pair vector.  Its exterior translation and similarity laws,
including the full derivative reflection of $V$, are
\begin{equation}
 \partial_{z^a}B_\nu=-p_a^-B_\nu,\qquad
 (z\cdot\partial+1)B_\nu=-HB_\nu.
 \label{P:exterior-laws}
\end{equation}
Translation fixes the insertion at infinity; reflection changes its
generator to the lowering generator with the indicated sign.  For
the second formula the exterior vector contributes weight $-6$ and
the reflected interior pair dimension five, giving total degree
$-1-H$.  These statements need only similarity covariance of $V$.

For a smooth Maxwell array, set
$Q_{\nu ab}(z)=\ip{B_\nu(z)}{F_{ab}}$.
Equations~\eqref{P:Maxwell-null-vectors} and
\eqref{P:exterior-laws} make $Q$ homogeneous of degree $-3$ and
make it and its Hodge dual in $a,b$ divergence-free.  Rotation
covariance and antisymmetry leave precisely
\begin{equation}
 Q_{\nu ab}=A\,|z|^{-4}(\delta_{\nu a}z_b-\delta_{\nu b}z_a)
            +B\,|z|^{-4}\epsilon_{\nu abd}z_d.
 \label{P:Maxwell-tensors}
\end{equation}
For $n=z/|z|$, the divergence of the first tensor is
$|z|^{-4}(\delta_{\nu b}-4n_\nu n_b)$, while that of the second is
zero.  Hence $A=0$.  Hodge duality interchanges these tensors up to
orientation sign, so dual divergence gives $B=0$.  Dense smooth
arrays in the energy-two rotation types and continuity of the fixed
exterior pairing prove the general assertion.
\end{proof}

\subsection{Harmonic removal of the residual}

Consider the short Euclidean scalar function
\begin{equation}
 G_{\mu\nu}(x;y,z)
   =W(\ell_\infty\ell(x)V_\mu(y)V_\nu(z)),
 \qquad x,y,z\text{ distinct}.
 \label{P:G-definition}
\end{equation}
By \eqref{P:accessible-decoupling} and the leg identities,
\begin{equation}
 \partial_x^2G=0,\qquad
 \partial_{x^a}G_{\mu\nu}
       =-W(\ell_\infty J_a(x)V_\mu(y)V_\nu(z)).
 \label{P:G-harmonic}
\end{equation}
For the second equation use
$\partial\ell=V-J-\partial S$; the accessible terms vanish.
For the first, $\partial^2\ell=\Theta-\partial^2S$, whose two
terms likewise vanish in this correlation.

For fixed $y,z$, $G$ is bounded as $|x|\to\infty$.
To see the quantitative bound, transport the two scalars colliding
at infinity by a proper spherical inversion, adjusted by a reflection,
to a finite radial ball, placing the vector pair exterior.  The
separation $r'$ of the inner scalar pair is comparable to $|x|^{-1}$.
Its norm is at most $C(r')^{-4}$ by
\eqref{P:radial-contraction} and $H\geq0$, starting from a fixed
inner radius uniformly in angle.  The scalar inversion factor of the
moving point is $O(|x|^{-4})$, and cancels this bound.  The fixed
exterior pair has finite norm, proving boundedness.

We next determine the possible singularities at the two finite
points $y$ and $z$.  Harmonicity gives two radial powers in each
spherical harmonic.  Pair-space positivity will identify a singular
power with an actual low-energy spectral projection; the preceding
representation and flux results can then test that projection.

Project the harmonic function onto a spherical harmonic of degree
$q$ about $y$.  With $r=|x-y|$, its radial coefficient has the form
\begin{equation}
 g_q(r)=A_qr^q+D_qr^{-q-2}.
 \label{P:harmonic-coefficient}
\end{equation}
Choose a small radial ball centered at $y$, let $v_q$ be the harmonic
projection of the inner pair ket
$|\ell(y+r_0n)V_\mu(y)\rangle$ at fixed $r_0$, and let $B$ be the
fixed exterior ket.  The contraction identity gives
\begin{equation}
 r^5g_q(r)=r_0^5\ip{B}{(r/r_0)^Hv_q}.
 \label{P:harmonic-spectral}
\end{equation}
The spectral pairing
$\mu(E)=\ip{B}{E_H(E)v_q}$ is a complex measure of finite total
variation, at most $\norm{B}\norm{v_q}$.  Thus the right side is a
Laplace transform of a finite measure on $[0,\infty)$.
The two powers on the left have exponents $q+5$ and $3-q$.
If $q>3$, boundedness of this transform at $r=0$ kills $D_q$.
Otherwise uniqueness of finite-measure Laplace transforms identifies
$D_q$ with the actual point spectral projection at energy $3-q$.
One can verify uniqueness by setting $u=e^{-h}$: all integer
Laplace parameters give the polynomial moments on $[0,1]$, which
determine the finite measure.  In particular, a continuous spectral
component cannot imitate the singular power in
\eqref{P:harmonic-coefficient}.

The projected array transforms in $\mathcal Y_q\otimes\mathbf4$,
where $\mathcal Y_q$ is the degree-$q$ harmonic representation.
At $q=3$ the possible energy is zero, and at $q=2$ it is one.  Only
scalars occur there by Lemma~\ref{P:radial-low-energy}, whereas
$\mathcal Y_q\otimes\mathbf4$ contains a scalar only for $q=1$.
Both coefficients therefore vanish.  At $q=1$ the energy is two;
the possible types are the scalar and Maxwell two-forms.  The latter
have zero exterior overlap by Lemma~\ref{P:Maxwell-exterior}.  The
unique scalar intertwiner leaves precisely a dipole parallel to the
vector index.  The remaining $q=0$ monopole is killed by integrating
\eqref{P:G-harmonic} over a small sphere about $y$:
$F_J(V)=0$, whereas the radial derivative of $D_0r^{-2}$ has flux
$-2|S^3|D_0$.  Repeat at $z$.

Convergent harmonic expansions on punctured-ball annuli now show
that subtraction of these two dipoles removes both singularities.
The remainder is entire and bounded at infinity, hence constant in
$x$.  We have proved
\begin{equation}
 G_{\mu\nu}(x;y,z)=
 D_\nu(y,z)\frac{(x-y)_\mu}{|x-y|^4}
 +\widetilde D_\mu(y,z)\frac{(x-z)_\nu}{|x-z|^4}
 +C_{\mu\nu}(y,z).
 \label{P:G-decomposition}
\end{equation}
The constant is the limit at $x$ infinity, so its similarity laws are
immediate: it depends only on $y-z$ and is homogeneous of degree $-6$.

The only possible finite singularities are now the two displayed
dipoles.  A tangential polarization kills them simultaneously after
placing the two vector points at inverse radial positions.  This
choice turns the bounded remainder into a bound on a genuine pair
norm, which excludes its low radial energies.

Take radial center $x=0$, choose $|u|=1$, put $y=u/r$, $z=ru$, and
contract both vector indices with a fixed $w\perp u$.  Both dipoles
in \eqref{P:G-decomposition} vanish, and its last term is $O(r^6)$.
For clarity, the full radial reflection law for $V=J+\partial L$ is
\begin{equation}
 V_\mu(x)^\dagger
 =|x|^{-6}R_{\mu\nu}(x)V_\nu(Ix)
      -4x_\mu|x|^{-6}L(Ix),\qquad
 R_{\mu\nu}=\delta_{\mu\nu}-2\widehat x_\mu\widehat x_\nu.
 \label{P:V-reflection}
\end{equation}
The second term is the derivative of $|x|^{-4}$ in the scalar
reflection formula.  Tangential contraction removes it and leaves
$Rw=w$.  Radial splitting therefore gives
\begin{equation}
 \big\|\,|\ell(0)V_w(ru)\rangle\,\big\|^2
     =r^{-6}G_{ww}(0;u/r,ru)=O(1).
 \label{P:tangential-norm}
\end{equation}
This bounds an actual Hilbert norm.  If $\psi$ is the corresponding
ket at a fixed radius $r_0$ and $s=r/r_0$, contraction gives
\begin{equation}
 \big\|\,|\ell(0)V_w(ru)\rangle\,\big\|^2
    =s^{-10}\int_{[0,\infty)}s^{2h}\,
                   \dd\ip{\psi}{E_H(h)\psi}.
 \label{P:no-low-energy}
\end{equation}
A nonzero spectral mass on $[0,5-\varepsilon]$ would make this at
least a positive constant times $s^{-2\varepsilon}$, contrary to
\eqref{P:tangential-norm}.  Taking all $\varepsilon>0$ excludes the
whole spectrum below five, in particular overlap with
$|J_\rho\rangle$ at energy three.

Use \eqref{P:ell-J-L} and its conformal transports to remove the
$\partial L$ part of $V$ in that overlap.  The remaining three-point
function has the similarity form
\begin{equation}
 \langle J_\rho|\ell(0)J_\mu(z)\rangle
      =|z|^{-2}(a\delta_{\rho\mu}
                      +b\widehat z_\rho\widehat z_\mu).
 \label{P:ell-JJ-form}
\end{equation}
Tangential testing forces $a=0$.  Its divergence is
$|z|^{-3}(b-2a)\widehat z_\rho$, so conservation forces $b=0$.
Primary conformal transport yields
\begin{equation}
 W(\ell_\infty JJ)=0.
 \label{P:ell-JJ}
\end{equation}
Expand the already zero $W(\ell_\infty VV)$ using $V=J+\partial L$.
Equations~\eqref{P:ell-J-L} and \eqref{P:ell-JJ} leave the mixed
second derivative of $W(\ell_\infty L(y)L(z))$ equal to zero.
Translation covariance makes this a function of $y-z$ with zero
Hessian, and its degree $-2$ excludes a nonzero affine function.
Consequently
\begin{equation}
 W(\ell_\infty LL)=0.
 \label{P:ell-LL}
\end{equation}

The residual already decouples from the $JL$, $JJ$ and $LL$
three-point functions.  To test its own scalar norm, we need one more
identity involving the stress tensor.  We first remove the allowed
$Jt$ tensor structure with a centered special flux, and then use a
centered dilation flux on the remaining $Lt$ function.

It remains to use the stress flux.  Write $R=|z-y|$ and
$n=(z-y)/R$.  Rotation and scale covariance, symmetry in
$\rho,\sigma$, and the absence of a nonzero parity-odd symmetric
tensor give the general form
\begin{align}
 W(\ell_\infty J_\mu(y)t_{\rho\sigma}(z))
  =R^{-5}\big[&A n_\mu\delta_{\rho\sigma}
       +B(n_\rho\delta_{\sigma\mu}+n_\sigma\delta_{\rho\mu})
       +D n_\mu n_\rho n_\sigma\big].
 \label{P:ell-Jt-general}
\end{align}
Trace gives $4A+2B+D=0$.  Vector divergence gives the coefficients
$-2A+2B$ and $-2D-12B$ of, respectively,
$R^{-6}\delta_{\rho\sigma}$ and $R^{-6}n_\rho n_\sigma$.
Thus $A=B=\gamma$ and $D=-6\gamma$, giving
\begin{equation}
 W(\ell_\infty J_\mu(y)t_{\rho\sigma}(z))
 =\gamma R^{-5}
 (n_\mu\delta_{\rho\sigma}+n_\rho\delta_{\sigma\mu}
       +n_\sigma\delta_{\rho\mu}-6n_\mu n_\rho n_\sigma).
 \label{P:ell-Jt}
\end{equation}
This also obeys tensor divergence.  Take the special vector field
centered at $y$,
$X_\alpha^\sigma=2(z-y)_\alpha(z-y)^\sigma-R^2\delta_\alpha{}^\sigma$.
The contraction of its angular part with the tensor in
\eqref{P:ell-Jt} is
\begin{equation}
 n^\rho(2n_\alpha n^\sigma-\delta_\alpha{}^\sigma)
 (n_\mu\delta_{\rho\sigma}+n_\rho\delta_{\sigma\mu}
       +n_\sigma\delta_{\rho\mu}-6n_\mu n_\rho n_\sigma)
 =-\delta_{\alpha\mu}-2n_\alpha n_\mu.
 \label{P:special-angular}
\end{equation}
Its sphere average is $-3\delta_{\alpha\mu}/2$, since
$\langle n_\alpha n_\mu\rangle=\delta_{\alpha\mu}/4$.
The field contributes $R^2$ and the area $R^3$, canceling $R^{-5}$.
The flux is therefore
$-3\gamma|S^3|\delta_{\alpha\mu}/2$.  It is zero by
\eqref{P:conformal-flux}, because a centered special transformation
annihilates the primary $J$.  Hence $\gamma=0$.

\Needspace{6\baselineskip}
\begin{theorem}[Return of the scalar wave]
\label{P:return}
The formal scalar vacuum wave equals the accessible scalar:
\begin{equation}
 l=S,\qquad \ell=0.
 \label{P:return-equation}
\end{equation}
\end{theorem}

\begin{proof}
Expand $W(\ell_\infty V(y)T(z))=0$, which is a case of
\eqref{P:accessible-decoupling}, using
\eqref{P:euclidean-identities}.  The $Jt$ term has just vanished;
the terms with derivatives of $L$ vanish by
\eqref{P:ell-J-L} and \eqref{P:ell-LL}.  What remains is
$\partial_yW(\ell_\infty L(y)t(z))=0$.
Translation covariance and homogeneity of degree $-4$ force that
function to be zero.  Its centered dilation flux around $L$,
using \eqref{P:conformal-flux}, now gives
\begin{equation}
 0=2c\,W(\ell_\infty L).
 \label{P:final-flux}
\end{equation}
The energy-slice calculation in Proposition~\ref{E:accessible-scalar}
gave $D(e)=B(e,e)=b$ after excluding the massless contribution.  Thus
the constant two-point density of $S$ with $l-S$ vanishes: the scalar
wave channels are orthogonal, not merely their labels on one time
slice.  It follows that $W(\ell_\infty L)$ equals the
scalar two-point coefficient $W(\ell_\infty\ell)$, which is the
nonnegative radial norm $\norm{|\ell\rangle}^2$ and is strictly
positive for a nonzero scalar wave.  Since $c\ne0$,
\eqref{P:final-flux} makes that norm zero.  Therefore $l=S$.
\end{proof}
